\begin{abstract}
Le séquençage d'ADN rapide rend possible de nouvelles applications en biologie et en médecine.
Néanmoins, l'analyse d'un volume de plus en plus grand de données génomiques requiert des ressources
informatiques et un temps de calcul massifs, d'où l'importance de développer des outils d'analyse
efficaces.

La métagénomique étudie le matériel génétique prélevé dans les microbiomes, des communautés de
microorganismes qui vivent dans des habitats très variés comme le sol, l'eau, le système digestif
humain et d'autres environnements. Une meilleure compréhension des microbiomes est nécessaire pour
améliorer notre capacité à les caractériser, ce qui aiderait à détecter les pathogènes plus
efficacement ou mieux caractériser les sols en agriculture, par exemple. Une méthode répandue pour
analyser des métagénomes consiste à comparer des séquences d'ADN prélevées dans un environnement
avec des bases de données contenant des génomes de référence. Cette approche est toutefois lente et
limitée par la qualité des bases de données. Une technique plus récente et plus flexible consiste à
entraîner des réseaux neuronaux pour identifier les espèces présentes dans un métagénome, mais elle
requiert un temps d'entraînement important et comporte des difficultés en lien avec l'interprétation
des résultats.

Ce mémoire présente de nouvelles techniques de classification taxonomique des séquences d'ADN
métagénomique, ce qui consiste à catégoriser les organismes dans des groupes génétiquement
apparentés. Les techniques sont basées sur le modèle Mamba, un modèle à espace d'états, ainsi que
des variantes modifiées par un traitement parallèle en flux (Mamba-Dual), un pré-traitement par CNN
(Mamba-CNN), et une combinaison des deux améliorations (Mamba-Dual-CNN). Sur des ensembles de
séquences d'ADN génétiquement diversifiées, Mamba-Dual-CNN classifie les données avec une précision
moyenne de $89,97\%$, ce qui se compare favorablement à un transformateur de taille similaire
($72,46\%$), le modèle BERTax ($86,88\%$) et les outils MMSeqs 2 ($88,20\%$) ainsi que Kraken 2
($75,51\%$). La durée d'entraînement des modèles à espace d'états testés est outre moins longue
qu'un transformateur de taille similaire, avec $0,164\;s$ pour un pas avec Mamba-Dual-CNN contre
$0,188\;s$ pour le transformateur.

\end{abstract}
